% sections/zh/1-introduction.tex -- 漏斗五段式；与 sections/en/1-introduction.tex 逐段对应。
\section{引言}\label{sec:intro}

多重度柔性作业车间调度问题描述这样的车间：工件只有少数几种类型，每种类型有多个副本，副本共享同一条工艺路线、每道工序同一个候选机器集与同一组加工时长。电子装配、注塑与批量机加工车间都是这种形态：一个排产窗口内有几千道工序，但只有十来种产品。\pending{计划员每个窗口要出一版可执行的调度；用 job-indexed 遗传算法时，副本数翻倍后求解质量明显下滑。}难点在结构而非规模：同类副本可互换使 job-indexed 表示存在大量等价解，连续配额取整为整数带来的损失与固定配额下排序带来的损失来源不同。

\pending{多重度 FJSP 的流体松弛能给出有意义的分流比例与下界，流体辅助的进化与强化学习方法近年连续出现。}这些方法把流体解当作初始化、状态特征或控制信号，没有说明离散调度的困难里哪一部分已被流体模型解释；学习组件因此面对整个改进问题，在分布外失效时也没有非学习候选兜底。

本文从一个精确恒等式出发：流体下界到可行 makespan 的差距等于只依赖整数配额的取整损失与在该配额下由排序决定的排序残差之和。该恒等式给出了一种分工：进化搜索在维度由类型结构决定的配额染色体上进行，学习器只学习排序残差，精确解码是接受与否的唯一裁决者。

\draftnote{以下贡献列表为预期贡献。按反推式写作纪律，终稿的贡献列表只能由 \texttt{docs/claim-verdicts.md} 中裁决为“成立”的 claim 反推得到；“部分成立”的须收窄到证据边界内，“不成立”的必须删除。}

本文的贡献如下。
\begin{enumerate}
  \item 配额压缩表示及其流体-配额几何（第 \ref{sec:quota} 节）：基因型维度不随副本数增长，配额迁移保持各类型总量守恒，largest-remainder 配额的取整损失有与多重度无关的上界（命题 \ref{prop:rounding}）。
  \item 残差限定的商图排序器（第 \ref{sec:ranker} 节）：只预测排序残差，标签按流体下界归一化，输出不低于解析下界。
  \item 保底预选（第 \ref{sec:safeguard} 节）：学习器、流体启发式与随机三源并集在精确解码下裁决，局部不劣于单用启发式（命题 \ref{prop:safeguard}），并以可证伪的解码次数门槛评判。
  \item \pending{在 59 个固定设计算例、Brandimarte 公开基准与多重度外推档上的实验研究，含取整差距对其理论界的因子效应分析。}
\end{enumerate}

本文其余部分组织如下。第 \ref{sec:related} 节在流体辅助与学习辅助调度中定位本文；第 \ref{sec:problem} 节给出问题、流体松弛与差距恒等式；第 \ref{sec:method} 节给出方法及其结构性质；第 \ref{sec:experiments} 节报告实验；第 \ref{sec:conclusion} 节总结。
